% Archived detailed essential-spectrum argument
%
% This file preserves the detailed coupled/decoupled resolvent route that was
% formerly included in theory.tex and theory-zh.tex.  It is retained as a
% technical proof attempt and reference record, but it is no longer the main
% argument used in the draft to determine numerical scan intervals.  This
% fragment is intentionally not \input by either main document.

\section*{English archived argument}

\begin{proposition}[Essential spectrum for two different half-guides]
\label{prop:asymmetric-essential-spectrum}
Assume that $\rho$ agrees with $\rho^-$ on $W^-$ and with $\rho^+$ on $W^+$,
while the region where it differs from these periodic backgrounds is contained
in the bounded center cell $\overline{\mathcal C^0}$.  Then
\begin{equation}
  \boxed{\quad
  \sigma_{\mathrm{ess}}(A(\beta))
  =\Sigma_{\mathrm{ess}}(\beta)
  :=\sigma(A^-(\beta))\cup\sigma(A^+(\beta)).\quad}
  \label{eq:projected-spectrum}
\end{equation}
No equality between the two background spectra is required.
\end{proposition}
\begin{proof}
\begin{enumerate}[label=\textit{Step \arabic*.},leftmargin=*]
\item \emph{Construct both realizations on one Hilbert space and locate a
common resolvent point.}
The bounds in \eqref{eq:material} give the common Hilbert space
\[
  H:=L^2(B,\rho\,\dd x\dd y),
  \qquad
  \langle u,v\rangle_H
  :=\int_B\rho\,u\overline v\,\dd x\dd y.
\]
Here
$0<\rho_{\min}:=\rho_-\leq\rho\leq\rho_+=:\rho_{\max}<\infty$.
As a set, $H=L^2(B)$, and its weighted norm is equivalent to the ordinary
$L^2$ norm.  In particular, $H^1(B)\subset H$.  Since the ports have measure
zero, restriction to the three open pieces gives the orthogonal decomposition
\[
  H=H^-\oplus H^0\oplus H^+,
  \qquad
  \begin{aligned}
    H^-&:=L^2(W^-,\rho^-\,\dd x\dd y),\\
    H^0&:=L^2(\mathcal C^0,\rho\,\dd x\dd y),\\
    H^+&:=L^2(W^+,\rho^+\,\dd x\dd y).
  \end{aligned}
\]

For any one of the domains used below, put
\[
  H^1(\Delta,D):=\{u\in H^1(D):\Delta u\in L^2(D)\}.
\]
The coupled realization is the operator
\[
  A(\beta):D(A(\beta))\subset H\longrightarrow H,
  \qquad A(\beta)u=-\frac1\rho\Delta u,
\]
with
\[
  D(A(\beta))=\left\{u\in H^1(\Delta,B):
  \begin{array}{l}
    \left.u\right|_{y=d}
      =e^{\ii\beta d}\left.u\right|_{y=0},\\[-1mm]
    \left.\partial_yu\right|_{y=d}
      =e^{\ii\beta d}\left.\partial_yu\right|_{y=0}
  \end{array}\right\}.
\]
The value equality is an $H^{1/2}$ trace equality and the derivative equality
is understood in the weak $H^{-1/2}$ trace space.  To realize this operator by
a closed form, set
\[
  V_\beta:=H^1_\beta(B),
  \qquad
  a_\beta(u,v):=\int_B\nabla u\cdot\nabla\overline v\,\dd x\dd y,
  \qquad u,v\in V_\beta.
\]
The space of functions $e^{\ii\beta y}\phi(x,y)$ with $\phi$ smooth,
$d$-periodic in $y$, and compactly supported in $x$ is contained in
$V_\beta$ and is dense in $H$: this follows from the unitary gauge
$u\mapsto e^{-\ii\beta y}u$, periodic mollification in $y$, and cutoffs in
$x$.  Thus $a_\beta$ is densely defined.  It is sesquilinear, symmetric, and
nonnegative.  Moreover,
\[
  \|u\|_{a_\beta}^2
  :=a_\beta(u,u)+\|u\|_H^2
\]
is equivalent to the $H^1(B)$ norm by the bounds on $\rho$.  The Sobolev trace
theorem makes $V_\beta$ a closed subspace of $H^1(B)$, so this form norm is
complete and $a_\beta$ is closed; see also
\cite[Chapter~3]{McLean2000}.  The first representation theorem for closed,
densely defined, nonnegative forms
\cite[Chapter~VI, Section~2]{Kato1995} therefore produces a self-adjoint
nonnegative operator in $H$.  Green's identity and cancellation of the two
quasiperiodic horizontal boundary terms identify that associated operator
with the differential realization $A(\beta)$ displayed above; the domain
identification is detailed in the remark following this step.

For the decoupled realization, impose homogeneous Dirichlet conditions at the
ports and define
\[
  A_D^- (\beta):D(A_D^-(\beta))\subset H^-\longrightarrow H^- ,
  \qquad A_D^-(\beta)u^-=-\frac1{\rho^-}\Delta u^- ,
\]
\[
  D(A_D^-(\beta)):=\left\{u^-\in H^1(\Delta,W^-):
  \begin{array}{l}
    \left.u^-\right|_{\Gamma^-}=0,\\
    \left.u^-\right|_{y=d}
      =e^{\ii\beta d}\left.u^-\right|_{y=0},\\[-1mm]
    \left.\partial_yu^-\right|_{y=d}
      =e^{\ii\beta d}\left.\partial_yu^-\right|_{y=0}
  \end{array}\right\},
\]
\[
  A_D^0(\beta):D(A_D^0(\beta))\subset H^0\longrightarrow H^0,
  \qquad A_D^0(\beta)u^0=-\frac1\rho\Delta u^0,
\]
\[
  D(A_D^0(\beta)):=\left\{u^0\in H^1(\Delta,\mathcal C^0):
  \begin{array}{l}
    \left.u^0\right|_{\Gamma^-}
      =\left.u^0\right|_{\Gamma^+}=0,\\
    \left.u^0\right|_{y=d}
      =e^{\ii\beta d}\left.u^0\right|_{y=0},\\[-1mm]
    \left.\partial_yu^0\right|_{y=d}
      =e^{\ii\beta d}\left.\partial_yu^0\right|_{y=0}
  \end{array}\right\},
\]
and
\[
  A_D^+(\beta):D(A_D^+(\beta))\subset H^+\longrightarrow H^+,
  \qquad A_D^+(\beta)u^+=-\frac1{\rho^+}\Delta u^+,
\]
\[
  D(A_D^+(\beta)):=\left\{u^+\in H^1(\Delta,W^+):
  \begin{array}{l}
    \left.u^+\right|_{\Gamma^+}=0,\\
    \left.u^+\right|_{y=d}
      =e^{\ii\beta d}\left.u^+\right|_{y=0},\\[-1mm]
    \left.\partial_yu^+\right|_{y=d}
      =e^{\ii\beta d}\left.\partial_yu^+\right|_{y=0}
  \end{array}\right\}.
\]
All these boundary equalities have the same weak trace meaning as above.  Their
form domains are
\[
\begin{aligned}
  V_D^-&:=\{u^-\in H^1_\beta(W^-):
    \left.u^-\right|_{\Gamma^-}=0\},\\
  V_D^0&:=\{u^0\in H^1_\beta(\mathcal C^0):
    \left.u^0\right|_{\Gamma^-}
    =\left.u^0\right|_{\Gamma^+}=0\},\\
  V_D^+&:=\{u^+\in H^1_\beta(W^+):
    \left.u^+\right|_{\Gamma^+}=0\}.
\end{aligned}
\]
On $V_{\mathrm{dec}}:=V_D^-\oplus V_D^0\oplus V_D^+$ define
\[
  a_{\mathrm{dec}}(U,V)
  :=\sum_{\star\in\{-,0,+\}}
    \int_{D^\star}\nabla u^\star\cdot\nabla\overline{v^\star}
      \,\dd x\dd y,
  \qquad
  D^-:=W^-,\quad D^0:=\mathcal C^0,\quad D^+:=W^+.
\]
Smooth functions compactly supported in each open piece show that this form is
dense in $H$.  It is sesquilinear, symmetric, and nonnegative.  Its form norm
is equivalent to the sum of the three $H^1(D^\star)$ norms.  Each local form
domain is closed in the corresponding $H^1$ space by continuity of the trace,
so $V_{\mathrm{dec}}$ is complete in the form norm and
$a_{\mathrm{dec}}$ is closed.  The same representation theorem, followed by
Green's identity on each piece, therefore identifies its associated operator
as the self-adjoint, nonnegative direct sum
\[
  A_{\mathrm{dec}}(\beta)
  :=A_D^-(\beta)\oplus A_D^0(\beta)\oplus A_D^+(\beta)
  \quad\text{in }H.
\]
The compact embedding $H^1(\mathcal C^0)\hookrightarrow H^0$ also implies
that $A_D^0(\beta)$ has compact resolvent
\cite[Chapter~3]{McLean2000}.

We have now proved that $A(\beta)$ and $A_{\mathrm{dec}}(\beta)$ are
self-adjoint, nonnegative, and hence closed operators on the same Hilbert
space.  Their spectra lie in $[0,\infty)$
\cite[Chapter~V]{Kato1995}, and consequently
\[
  -1\in\rho(A(\beta))\cap\rho(A_{\mathrm{dec}}(\beta)).
\]
Here $\rho(T)$ denotes the resolvent set of an operator $T$.
In particular, both resolvents are everywhere defined bounded operators,
\[
  (A(\beta)+I)^{-1},\ (A_{\mathrm{dec}}(\beta)+I)^{-1}
  \in\mathcal B(H),
  \qquad
  \|(A(\beta)+I)^{-1}\|,\
  \|(A_{\mathrm{dec}}(\beta)+I)^{-1}\|\leq1,
\]
and their common domain as bounded operators is all of $H$.  The same
conclusion can be checked directly by Lax--Milgram: for every $f\in H$, the
forms
\[
  a_\beta(u,v)+\langle u,v\rangle_H,
  \qquad
  a_{\mathrm{dec}}(U,V)+\langle U,V\rangle_H
\]
are continuous and coercive in their respective form norms, and therefore the
two variational equations with right-hand side $\langle f,\cdot\rangle_H$
have unique solutions and satisfy the displayed norm bounds
\cite[Lemma~2.32]{McLean2000}.  The closed-form theorem supplies the
self-adjoint realizations, while Lax--Milgram independently verifies the
common resolvent point and its estimate; no spectral-invariance conclusion is
used in this argument.

The resolvent difference
\[
  (A(\beta)+I)^{-1}-(A_{\mathrm{dec}}(\beta)+I)^{-1}
\]
is therefore a well-defined member of $\mathcal B(H)$.  The Weyl resolvent
form of the essential-spectrum theorem states that, for two self-adjoint
operators on the same Hilbert space and a common resolvent point $z$, a
compact resolvent difference at $z$ implies equality of their essential
spectra \cite[Theorem~XIII.14]{ReedSimon1978}.  At $z=-1$ all of its
hypotheses except that compactness have now been verified.  The only remaining
hypothesis needed for the Weyl-type invariance argument is to prove that the
above resolvent difference is compact on $H$.
\end{enumerate}

\begin{proofremark}[Form domains and operator domains]
The space $V_\beta$ is the form domain; it is not the operator domain.  If
$T_\beta$ denotes the operator associated with $a_\beta$, then abstractly
\[
\begin{aligned}
  D(T_\beta)=\bigl\{u\in V_\beta:{}&\ \text{there exists }f\in H
  \text{ such that}\\[-1mm]
  &a_\beta(u,v)=\langle f,v\rangle_H
  \text{ for every }v\in V_\beta\bigr\},
  \qquad T_\beta u=f.
\end{aligned}
\]
This is exactly the $H^1(\Delta,B)$ trace domain used for $A(\beta)$ above,
and hence agrees with the realization in Fliss~\cite[Section~3]{Fliss2013}.
Indeed, one implication follows from the weak Green identity: the horizontal
boundary terms cancel because both the value and weak $\partial_y$ traces are
$\beta$-quasiperiodic.  Conversely, testing the abstract identity first
with functions compactly supported in the interior of $B$ gives
$-\Delta u=\rho f$ in distributions, so $\Delta u\in L^2(B)$.  Applying the
weak Green identity to arbitrary $\beta$-quasiperiodic test functions then
forces
\[
  \left.\partial_yu\right|_{y=d}
  =e^{\ii\beta d}\left.\partial_yu\right|_{y=0}
\]
in $H^{-1/2}$; the value quasiperiodicity was already imposed by
$u\in V_\beta$.  The local Dirichlet realizations are identified in the same
way: zero value traces at the artificial ports are imposed in
$V_D^\star$, while Green's identity recovers the remaining natural
quasiperiodic derivative condition.  The weak Green identities and normal
traces used here are those for $H^1(\Delta,D)$; see
\cite[Chapters~3--4]{McLean2000}.
\end{proofremark}

\begin{enumerate}[label=\textit{Step \arabic*.},leftmargin=*,start=2]

\item \emph{Compare the coupled and decoupled resolvents.}
Both $A(\beta)$ and $A_{\mathrm{dec}}(\beta)$ are nonnegative, so
$-1$ belongs to their common resolvent set.  Cutting or restoring transmission
across $\Gamma^-\cup\Gamma^+$ changes the resolvent by a compact operator:
\begin{equation*}
  (A(\beta)+1)^{-1}-(A_{\mathrm{dec}}(\beta)+1)^{-1}
  \quad\text{is compact on }H.
\end{equation*}
Indeed, localized resolvents have additional elliptic regularity near the
ports.  The map from a volume right-hand side to its port Dirichlet traces is
therefore bounded into a Sobolev space compactly embedded in $H^{1/2}$, while
the Poisson solution map from $H^{1/2}$ port data into $H$ is bounded.  The
resolvent difference factors through these two maps and is compact.
Weyl's essential spectrum theorem
\cite[Theorem~XIII.14]{ReedSimon1978} now yields
\[
  \sigma_{\mathrm{ess}}(A(\beta))
  =\sigma_{\mathrm{ess}}(A_{\mathrm{dec}}(\beta)).
\]

\item \emph{Compute the essential spectrum of the direct sum.}
The essential spectrum of a finite orthogonal direct sum is the union of the
essential spectra of its summands.  Since $A_D^0(\beta)$ has compact
resolvent, $\sigma_{\mathrm{ess}}(A_D^0(\beta))=\varnothing$, and hence
\[
  \sigma_{\mathrm{ess}}(A(\beta))
  =\sigma_{\mathrm{ess}}(A_D^-(\beta))
   \cup\sigma_{\mathrm{ess}}(A_D^+(\beta)).
\]

\item \emph{Identify each half-guide contribution.}
Fix $\star\in\{-,+\}$.  If
$\lambda\in\sigma(A^\star(\beta))$, Floquet decomposition produces Bloch
wave packets supported over an increasing number of cells whose residuals
tend to zero.  Translating these packets arbitrarily far into $W^\star$ and
cutting them off away from the finite port gives a singular Weyl sequence for
$A_D^\star(\beta)$.  Thus
\[
  \sigma(A^\star(\beta))
  \subset\sigma_{\mathrm{ess}}(A_D^\star(\beta)).
\]
Conversely, local compactness forces every singular Weyl sequence for
$A_D^\star(\beta)$ to escape from its finite port.  Cutting it off in cells
far from that port and extending it into the two-sided periodic guide gives a
Weyl sequence for $A^\star(\beta)$.  Consequently,
\[
  \sigma_{\mathrm{ess}}(A_D^\star(\beta))
  \subset\sigma(A^\star(\beta)),
\]
and therefore
\[
  \sigma_{\mathrm{ess}}(A_D^\star(\beta))
  =\sigma(A^\star(\beta)).
\]
The full spectrum of $A_D^\star(\beta)$ may additionally contain discrete
port-induced eigenvalues; they do not enter this essential-spectrum identity.
Applying the identity for $\star=-$ and $\star=+$ in Step~3 proves
\eqref{eq:projected-spectrum}.
\end{enumerate}
\end{proof}

\clearpage
\section*{中文归档论证}

本节保存原中文草稿中的详细预解式论证路线。它目前不再作为主草稿中决定数值扫描区间的主要论证，也不会由中文主文件通过 \verb|\input| 载入。

\begin{proposition}[两个不同半波导对应的本质谱]
\label{prop:asymmetric-essential-spectrum}
假设 $\rho$ 在 $W^-$ 上等于 $\rho^-$、在 $W^+$ 上等于 $\rho^+$，并且它偏离这两个周期背景的区域包含在有界中心胞元 $\overline{\mathcal C^0}$ 中。则
\begin{equation}
  \boxed{\quad
  \sigma_{\mathrm{ess}}(A(\beta))
  =\Sigma_{\mathrm{ess}}(\beta)
  :=\sigma(A^-(\beta))\cup\sigma(A^+(\beta)).\quad}
  \label{eq:projected-spectrum}
\end{equation}
这里不要求两个背景谱相等。
\end{proposition}
\begin{proof}
\begin{enumerate}[label=\textit{步骤 \arabic*.},leftmargin=*]
\item \emph{在同一 Hilbert 空间中构造两种实现，并确定一个公共预解点。}
由 \eqref{eq:material} 中的界可取共同的 Hilbert 空间
\[
  H:=L^2(B,\rho\,\dd x\dd y),
  \qquad
  \langle u,v\rangle_H
  :=\int_B\rho\,u\overline v\,\dd x\dd y.
\]
这里 $0<\rho_{\min}:=\rho_-\leq\rho\leq\rho_+=:\rho_{\max}<\infty$。
作为函数集合，$H=L^2(B)$，其加权范数与通常的 $L^2$ 范数等价。特别地，
$H^1(B)\subset H$。由于两个端口均为零测集，对三个开区域作限制便得到正交分解
\[
  H=H^-\oplus H^0\oplus H^+,
  \qquad
  \begin{aligned}
    H^-&:=L^2(W^-,\rho^-\,\dd x\dd y),\\
    H^0&:=L^2(\mathcal C^0,\rho\,\dd x\dd y),\\
    H^+&:=L^2(W^+,\rho^+\,\dd x\dd y).
  \end{aligned}
\]

对于下文出现的任一区域 $D$，记
\[
  H^1(\Delta,D):=\{u\in H^1(D):\Delta u\in L^2(D)\}.
\]
耦合实现是算子
\[
  A(\beta):D(A(\beta))\subset H\longrightarrow H,
  \qquad A(\beta)u=-\frac1\rho\Delta u,
\]
其中
\[
  D(A(\beta))=\left\{u\in H^1(\Delta,B):
  \begin{array}{l}
    \left.u\right|_{y=d}
      =e^{\ii\beta d}\left.u\right|_{y=0},\\[-1mm]
    \left.\partial_yu\right|_{y=d}
      =e^{\ii\beta d}\left.\partial_yu\right|_{y=0}
  \end{array}\right\}.
\]
函数值等式是 $H^{1/2}$ 中的迹等式，导数等式则在弱 $H^{-1/2}$ 迹空间中理解。
为通过闭型实现这一算子，令
\[
  V_\beta:=H^1_\beta(B),
  \qquad
  a_\beta(u,v):=\int_B\nabla u\cdot\nabla\overline v\,\dd x\dd y,
  \qquad u,v\in V_\beta.
\]
所有形如 $e^{\ii\beta y}\phi(x,y)$ 的函数组成的空间包含于 $V_\beta$，其中
$\phi$ 光滑、关于 $y$ 以 $d$ 为周期，并且关于 $x$ 具有紧支撑；这个空间在 $H$
中稠密。事实上，这可由酉规范变换 $u\mapsto e^{-\ii\beta y}u$、$y$ 方向的周期
磨光以及 $x$ 方向的截断得到。因此，$a_\beta$ 稠密定义。它是半双线性、对称且
非负的。此外，型范数
\[
  \|u\|_{a_\beta}^2
  :=a_\beta(u,u)+\|u\|_H^2
\]
由 $\rho$ 的上下界可知与 $H^1(B)$ 范数等价。Sobolev 迹定理表明 $V_\beta$ 是
$H^1(B)$ 的闭子空间，故它关于该型范数完备，从而 $a_\beta$ 是闭型；另见
\cite[Chapter~3]{McLean2000}。闭的、稠密定义的非负型的第一表示定理
\cite[Chapter~VI, Section~2]{Kato1995} 因而在 $H$ 中给出一个自伴非负算子。
Green 恒等式以及上下两条准周期水平边界项的抵消表明，该关联算子正是上面给出的
微分实现 $A(\beta)$；具体的定义域识别见本步骤之后的注。

对于切断实现，在两个端口施加齐次 Dirichlet 条件，并定义
\[
  A_D^- (\beta):D(A_D^-(\beta))\subset H^-\longrightarrow H^- ,
  \qquad A_D^-(\beta)u^-=-\frac1{\rho^-}\Delta u^- ,
\]
\[
  D(A_D^-(\beta)):=\left\{u^-\in H^1(\Delta,W^-):
  \begin{array}{l}
    \left.u^-\right|_{\Gamma^-}=0,\\
    \left.u^-\right|_{y=d}
      =e^{\ii\beta d}\left.u^-\right|_{y=0},\\[-1mm]
    \left.\partial_yu^-\right|_{y=d}
      =e^{\ii\beta d}\left.\partial_yu^-\right|_{y=0}
  \end{array}\right\},
\]
\[
  A_D^0(\beta):D(A_D^0(\beta))\subset H^0\longrightarrow H^0,
  \qquad A_D^0(\beta)u^0=-\frac1\rho\Delta u^0,
\]
\[
  D(A_D^0(\beta)):=\left\{u^0\in H^1(\Delta,\mathcal C^0):
  \begin{array}{l}
    \left.u^0\right|_{\Gamma^-}
      =\left.u^0\right|_{\Gamma^+}=0,\\
    \left.u^0\right|_{y=d}
      =e^{\ii\beta d}\left.u^0\right|_{y=0},\\[-1mm]
    \left.\partial_yu^0\right|_{y=d}
      =e^{\ii\beta d}\left.\partial_yu^0\right|_{y=0}
  \end{array}\right\},
\]
以及
\[
  A_D^+(\beta):D(A_D^+(\beta))\subset H^+\longrightarrow H^+,
  \qquad A_D^+(\beta)u^+=-\frac1{\rho^+}\Delta u^+,
\]
\[
  D(A_D^+(\beta)):=\left\{u^+\in H^1(\Delta,W^+):
  \begin{array}{l}
    \left.u^+\right|_{\Gamma^+}=0,\\
    \left.u^+\right|_{y=d}
      =e^{\ii\beta d}\left.u^+\right|_{y=0},\\[-1mm]
    \left.\partial_yu^+\right|_{y=d}
      =e^{\ii\beta d}\left.\partial_yu^+\right|_{y=0}
  \end{array}\right\}.
\]
以上所有边界等式均与前文具有相同的弱迹含义。相应的型定义域为
\[
\begin{aligned}
  V_D^-&:=\{u^-\in H^1_\beta(W^-):
    \left.u^-\right|_{\Gamma^-}=0\},\\
  V_D^0&:=\{u^0\in H^1_\beta(\mathcal C^0):
    \left.u^0\right|_{\Gamma^-}
    =\left.u^0\right|_{\Gamma^+}=0\},\\
  V_D^+&:=\{u^+\in H^1_\beta(W^+):
    \left.u^+\right|_{\Gamma^+}=0\}.
\end{aligned}
\]
在 $V_{\mathrm{dec}}:=V_D^-\oplus V_D^0\oplus V_D^+$ 上定义
\[
  a_{\mathrm{dec}}(U,V)
  :=\sum_{\star\in\{-,0,+\}}
    \int_{D^\star}\nabla u^\star\cdot\nabla\overline{v^\star}
      \,\dd x\dd y,
  \qquad
  D^-:=W^-,\quad D^0:=\mathcal C^0,\quad D^+:=W^+.
\]
每个开区域中具有紧支撑的光滑函数表明该型在 $H$ 中稠密定义。它是半双线性、
对称且非负的，其型范数与三个 $H^1(D^\star)$ 范数平方和的平方根等价。由迹映射
的连续性，每个局部型定义域都是相应 $H^1$ 空间的闭子空间，故
$V_{\mathrm{dec}}$ 关于型范数完备，$a_{\mathrm{dec}}$ 是闭型。再次应用表示定理，
并在每个区域上使用 Green 恒等式，即可把关联算子识别为 $H$ 中的自伴非负直和
\[
  A_{\mathrm{dec}}(\beta)
  :=A_D^-(\beta)\oplus A_D^0(\beta)\oplus A_D^+(\beta).
\]
紧嵌入 $H^1(\mathcal C^0)\hookrightarrow H^0$ 还表明
$A_D^0(\beta)$ 具有紧预解式 \cite[Chapter~3]{McLean2000}。

至此已经证明，$A(\beta)$ 与 $A_{\mathrm{dec}}(\beta)$ 是同一 Hilbert 空间上的
自伴、非负并因而闭的算子。它们的谱均包含于 $[0,\infty)$
\cite[Chapter~V]{Kato1995}，所以
\[
  -1\in\rho(A(\beta))\cap\rho(A_{\mathrm{dec}}(\beta)).
\]
这里 $\rho(T)$ 表示算子 $T$ 的预解集。
特别地，两个预解式都是处处定义的有界算子，并且
\[
  (A(\beta)+I)^{-1},\ (A_{\mathrm{dec}}(\beta)+I)^{-1}
  \in\mathcal B(H),
  \qquad
  \|(A(\beta)+I)^{-1}\|,\
  \|(A_{\mathrm{dec}}(\beta)+I)^{-1}\|\leq1;
\]
作为有界算子，它们的定义域都是整个 $H$。也可以用 Lax--Milgram 定理直接检验
这一结论：对每个 $f\in H$，两个型
\[
  a_\beta(u,v)+\langle u,v\rangle_H,
  \qquad
  a_{\mathrm{dec}}(U,V)+\langle U,V\rangle_H
\]
分别在各自的型范数下连续且强制，因此以 $\langle f,\cdot\rangle_H$ 为右端的两个
变分方程都有唯一解，并满足上面的范数估计
\cite[Lemma~2.32]{McLean2000}。闭型表示定理给出自伴实现，而 Lax--Milgram 定理
则独立核验公共预解点及其估计；这一论证没有使用任何谱不变性结论。

因此，预解式之差
\[
  (A(\beta)+I)^{-1}-(A_{\mathrm{dec}}(\beta)+I)^{-1}
\]
是 $\mathcal B(H)$ 中定义良好的元素。Weyl 本质谱定理的预解式形式指出：若两个
自伴算子作用于同一 Hilbert 空间，$z$ 是其公共预解点，并且它们在 $z$ 处的预解式
之差紧，则二者的本质谱相等
\cite[Theorem~XIII.14]{ReedSimon1978}。取 $z=-1$ 时，除紧性之外的所有假设均已
验证。为了应用上述 Weyl 型不变性论证，唯一尚需证明的假设是：上面的预解式之差
在 $H$ 上是紧算子。
\end{enumerate}

\begin{proofremark}[型定义域与算子定义域]
$V_\beta$ 是型定义域，而不是算子定义域。若 $T_\beta$ 表示与 $a_\beta$ 关联的
算子，则抽象地
\[
\begin{aligned}
  D(T_\beta)=\bigl\{u\in V_\beta:{}&\ \text{存在 }f\in H\text{，使得}\\[-1mm]
  &a_\beta(u,v)=\langle f,v\rangle_H
  \text{ 对每个 }v\in V_\beta\text{ 均成立}\bigr\},
  \qquad T_\beta u=f.
\end{aligned}
\]
这恰好是上面用于 $A(\beta)$ 的 $H^1(\Delta,B)$ 迹定义域，因而与
Fliss~\cite[Section~3]{Fliss2013} 中的实现一致。事实上，一个方向直接来自弱 Green
恒等式：由于函数值迹与弱 $\partial_y$ 迹都满足 $\beta$-准周期条件，两条水平边界上的
项相互抵消。反过来，先以支撑于 $B$ 内部的函数检验上述抽象等式，可在分布意义下
得到 $-\Delta u=\rho f$，故 $\Delta u\in L^2(B)$。再对任意 $\beta$-准周期检验
函数应用弱 Green 恒等式，便可在 $H^{-1/2}$ 中推出
\[
  \left.\partial_yu\right|_{y=d}
  =e^{\ii\beta d}\left.\partial_yu\right|_{y=0};
\]
函数值的准周期性已经由 $u\in V_\beta$ 给出。局部 Dirichlet 实现也可用相同方式
识别：人工端口上的零函数值迹由 $V_D^\star$ 直接施加，其余自然的准周期导数条件
则由 Green 恒等式恢复。这里使用的是 $H^1(\Delta,D)$ 上的弱 Green 恒等式和法向
迹；参见 \cite[Chapters~3--4]{McLean2000}。
\end{proofremark}

\begin{enumerate}[label=\textit{步骤 \arabic*.},leftmargin=*,start=2]

\item \emph{比较耦合算子与切断算子的预解式。}
$A(\beta)$ 和 $A_{\mathrm{dec}}(\beta)$ 都非负，故 $-1$ 属于二者的公共预解集。在 $\Gamma^-\cup\Gamma^+$ 上切断或恢复传输条件只使预解式改变一个紧算子：
\begin{equation*}
  (A(\beta)+1)^{-1}-(A_{\mathrm{dec}}(\beta)+1)^{-1}
  \quad\text{是 }H\text{ 上的紧算子。}
\end{equation*}
具体而言，局部化预解式在端口附近具有更高的椭圆正则性，因而从体右端项到端口 Dirichlet 迹的映射取值于一个紧嵌入 $H^{1/2}$ 的 Sobolev 空间；另一方面，从 $H^{1/2}$ 端口数据到 $H$ 中解的 Poisson 映射有界。预解式之差通过这两个映射分解，故为紧算子。现在应用 Weyl 本质谱定理 \cite[Theorem~XIII.14]{ReedSimon1978}，得到
\[
  \sigma_{\mathrm{ess}}(A(\beta))
  =\sigma_{\mathrm{ess}}(A_{\mathrm{dec}}(\beta)).
\]

\item \emph{计算直和算子的本质谱。}
有限正交直和的本质谱等于各分量本质谱之并。由于 $A_D^0(\beta)$ 具有紧预解式，$\sigma_{\mathrm{ess}}(A_D^0(\beta))=\varnothing$，所以
\[
  \sigma_{\mathrm{ess}}(A(\beta))
  =\sigma_{\mathrm{ess}}(A_D^-(\beta))
   \cup\sigma_{\mathrm{ess}}(A_D^+(\beta)).
\]

\item \emph{分别识别两个半波导的贡献。}
固定 $\star\in\{-,+\}$。若 $\lambda\in\sigma(A^\star(\beta))$，Floquet 分解给出支撑在越来越多胞元上的 Bloch 波包，且其残量趋于零。把这些波包平移到 $W^\star$ 的任意深处，并在远离有限端口处截断，即得到 $A_D^\star(\beta)$ 的奇异 Weyl 序列。因此
\[
  \sigma(A^\star(\beta))
  \subset\sigma_{\mathrm{ess}}(A_D^\star(\beta)).
\]
反过来，局部紧性迫使 $A_D^\star(\beta)$ 的每个奇异 Weyl 序列逃离其有限端口。在远离端口的胞元中将其截断，再延拓到双侧周期波导，就得到 $A^\star(\beta)$ 的 Weyl 序列。因此
\[
  \sigma_{\mathrm{ess}}(A_D^\star(\beta))
  \subset\sigma(A^\star(\beta)),
\]
从而
\[
  \sigma_{\mathrm{ess}}(A_D^\star(\beta))
  =\sigma(A^\star(\beta)).
\]
$A_D^\star(\beta)$ 的完整谱还可能含有人工端口诱导的离散特征值，但它们不进入上述本质谱恒等式。分别对 $\star=-$ 和 $\star=+$ 使用该恒等式，并代入步骤 3，即得 \eqref{eq:projected-spectrum}。
\end{enumerate}
\end{proof}

